\documentclass[a4paper,12pt]{article}

\usepackage[utf8]{inputenc}
\usepackage[T2A]{fontenc}
\usepackage[russian]{babel}
\usepackage{amsmath}
\usepackage{amssymb}
\usepackage{graphicx}
\usepackage{csvsimple}
\usepackage{float}
\usepackage{booktabs}
\usepackage{geometry}
\usepackage{csquotes}
\usepackage[strings]{underscore}
\geometry{margin=2cm}


\title{Лабораторная работа №1\\Методы одномерного поиска}
\author{Георгий Лаврентьев, Дмитрий Коваль - студенты группы М3232}
\date{\today}

\begin{document}
\maketitle

\section{Постановка задачи}
В работе реализованы и сравнены методы одномерной минимизации:
\begin{itemize}
    \item пассивный поиск;
    \item дихотомия;
    \item золотое сечение;
    \item Фибоначчи;
    \item параболы;
    \item библиотечный метод Брента (\texttt{minimize\_scalar(..., method='brent')}).
\end{itemize}

Исследованы 4 функции:
\begin{itemize}
    \item $f_1(x)=(x-2)^2+1$ --- ``хорошая'' унимодальная;
    \item $f_{2.1}(x)=\begin{cases}
        \delta^2,& |x-x_{\min}|\le \delta\\
        (x-x_{\min})^2,& \text{иначе}
    \end{cases}$, где $x_{\min}=1$, $\delta=0.8$ (плато);
    \item $f_{2.2}(x)=\begin{cases}
        (x-3)^2,& x<3\\
        (x-3)^4,& x\ge 3
    \end{cases}$ (асимметричная кривизна);
    \item $f_3(x)=(x-1.5)^2+\sin(5x)$ --- многомодальная.
\end{itemize}

Диапазон точностей: $\varepsilon \in \{10^{-1}, 10^{-2}, \dots, 10^{-8}\}$, рабочий интервал $[-2, 5]$.

\begin{figure}[H]
\centering
\includegraphics[width=\textwidth]{results/plots/all_functions.png}
\caption{Графики всех исследуемых функций}
\end{figure}

\section{Организация эксперимента}
Для каждой унимодальной функции и каждого метода собирались:
\begin{itemize}
    \item число итераций;
    \item число вычислений функции;
    \item история сужения интервала неопределённости.
\end{itemize}

Полные таблицы по всем $\varepsilon$ автоматически сохранены в \texttt{results/tables/*.csv}.
Ниже приведены компактные сводки для $\varepsilon=10^{-8}$ и графики зависимостей.

\section{Результаты: унимодальные функции}

\subsection{$f_1(x)=(x-2)^2+1$}
\begin{table}[H]
\centering
\small
\begin{tabular}{lrrrr}
\toprule
Метод & Итерации & Вычисления & $x^*$ & $f(x^*)$ \\
\midrule
Пассивный поиск & 1 & 100000 & 1.999980000 & 1.0000000004 \\
Дихотомия & 29 & 59 & 1.999999989 & 1.0000000000 \\
Золотое сечение & 43 & 45 & 2.000000010 & 1.0000000000 \\
Фибоначчи & 42 & 46 & 1.999999995 & 1.0000000000 \\
Парабола & 3 & 5 & 2.000000000 & 1.0000000000 \\
Брент (SciPy) & 5 & 8 & 2.000000000 & 1.0000000000 \\
\bottomrule
\end{tabular}
\caption{Сравнение методов на $f_1$ при $\varepsilon=10^{-8}$}
\end{table}

\begin{figure}[H]
\centering
\includegraphics[width=\textwidth]{results/plots/func_1__comparison.png}
\caption{Зависимость числа итераций/вычислений от точности для $f_1$}
\end{figure}

\begin{figure}[H]
\centering
\includegraphics[width=\textwidth]{results/plots/func_1__interval_dynamics.png}
\caption{Динамика интервала неопределённости для $f_1$}
\end{figure}

\textbf{Краткий вывод:} на гладкой квадратичной функции парабола и Брент дают минимальные затраты; пассивный поиск самый дорогой по числу вычислений. \\
Из графика видно, что интервальные методы (дихотомия, золотое сечение, Фибоначчи) демонстрируют логарифмическую зависимость числа итераций от точности $\varepsilon$, что соответствует теоретическим оценкам.  \\
Метод парабол показывает почти постоянное число итераций, так как для квадратичной функции аппроксимация параболой является точной уже на первом шаге.  \\
Пассивный поиск не зависит от поведения функции. Для метода Брента отсутствуют данные для построения графиков динамики интервала неопределённости, так как используемая реализация (\texttt{scipy.optimize.minimize\_scalar(..., method=’brent’)}) не предоставляет информации о промежуточных итерациях и границах интервала в процессе работы алгоритма.  

\subsection{$f_{2.1}$: функция с плато}
\begin{table}[H]
\centering
\small
\begin{tabular}{lrrrr}
\toprule
Метод & Итерации & Вычисления & $x^*$ & $f(x^*)$ \\
\midrule
Пассивный поиск & 1 & 100000 & 0.200052001 & 0.6400000000 \\
Дихотомия & 29 & 59 & 0.200000007 & 0.6400000000 \\
Золотое сечение & 43 & 45 & 1.799999999 & 0.6400000000 \\
Фибоначчи & 42 & 46 & 0.200000005 & 0.6400000000 \\
Парабола & 45 & 47 & 1.500000000 & 0.6400000000 \\
Брент (SciPy) & 37 & 40 & 0.856888666 & 0.6400000000 \\
\bottomrule
\end{tabular}
\caption{Сравнение методов на $f_{2.1}$ при $\varepsilon=10^{-8}$}
\end{table}

\begin{figure}[H]
\centering
\includegraphics[width=\textwidth]{results/plots/func_2.1__comparison.png}
\caption{Зависимость числа итераций/вычислений от точности для $f_{2.1}$}
\end{figure}

\begin{figure}[H]
\centering
\includegraphics[width=\textwidth]{results/plots/func_2.1__interval_dynamics.png}
\caption{Динамика интервала неопределённости для $f_{2.1}$}
\end{figure}

\textbf{Краткий вывод:} из-за плато минимум не единственный, поэтому разные методы сходятся к разным точкам с одинаковым значением $f=0.64$. \\
А метод парабол ведёт себя совсем необычно (видно на первом графике). Потому что на функциях с плато значения функции в рассматриваемых точках почти совпадают, из-за чего аппроксимирующая парабола определяется плохо и не даёт информативного шага. В результате метод фактически делает почти одинаковые (малые или неэффективные) шаги независимо от требуемой точности. Поэтому сходимость определяется не точностью $\varepsilon$, а особенностями функции, и число итераций остаётся практически одинаковым для всех значений $\varepsilon$.

\subsection{$f_{2.2}(x)=\begin{cases}(x-3)^2,& x<3\\(x-3)^4,& x\ge 3\end{cases}$: асимметричная унимодальная}
\begin{table}[H]
\centering
\small
\begin{tabular}{lrrrr}
\toprule
Метод & Итерации & Вычисления & $x^*$ & $f(x^*)$ \\
\midrule
Пассивный поиск & 1 & 100000 & 3.000010000 & $1.00\cdot 10^{-20}$ \\
Дихотомия & 29 & 59 & 3.000000007 & $2.86\cdot 10^{-33}$ \\
Золотое сечение & 43 & 45 & 3.000000001 & $2.36\cdot 10^{-36}$ \\
Фибоначчи & 42 & 46 & 3.000000004 & $1.54\cdot 10^{-34}$ \\
Парабола & 151 & 153 & 3.000000000 & $8.77\cdot 10^{-26}$ \\
Брент (SciPy) & 54 & 57 & 3.000000018 & $9.68\cdot 10^{-32}$ \\
\bottomrule
\end{tabular}
\caption{Сравнение методов на $f_{2.2}$ при $\varepsilon=10^{-8}$}
\end{table}

\begin{figure}[H]
\centering
\includegraphics[width=\textwidth]{results/plots/func_2.2__comparison.png}
\caption{Зависимость числа итераций/вычислений от точности для $f_{2.2}$}
\end{figure}

\begin{figure}[H]
\centering
\includegraphics[width=\textwidth]{results/plots/func_2.2__interval_dynamics.png}
\caption{Динамика интервала неопределённости для $f_{2.2}$}
\end{figure}

\textbf{Замечание: }На этой асимметричной функции метод парабол не теряет минимум, но работает заметно тяжелее остальных методов по числу итераций и вычислений. Метод Брента, как и интервальные методы, устойчиво уточняет точку минимума. Наглядно это видно по следующим таблицам сравнения методов Параболы и Брента на этой функции: \\

\csvautotabular{\detokenize{results/tables/func_2.2__Парабола.csv}} \\

\csvautotabular{\detokenize{results/tables/func_2.2__Брент_(SciPy).csv}}
\\ 
\\

\textbf{Краткий вывод:} метод парабол на этой функции заметно проигрывает по затратам, тогда как интервальные методы и Брент устойчиво находят минимум возле $x=3$.

\section{Многомодальная функция и разведка}
Для $f_3(x)=(x-1.5)^2+\sin(5x)$ получены следующие результаты при $\varepsilon=10^{-8}$:

\begin{table}[H]
\centering
\small
\begin{tabular}{lrrrr}
\toprule
Метод & Итерации & Вычисления & $x^*$ & $f(x^*)$ \\
\midrule
Пассивный поиск & 1 & 100000 & 0.984059841 & $-0.7122701655$ \\
Дихотомия & 29 & 59 & 0.984052380 & $-0.7122701662$ \\
Золотое сечение & 43 & 45 & 2.146777366 & $-0.5476338287$ \\
Фибоначчи & 42 & 46 & 2.146777366 & $-0.5476338287$ \\
Парабола & 13 & 15 & 0.984052374 & $-0.7122701662$ \\
Брент (SciPy) & 13 & 16 & 2.146777364 & $-0.5476338287$ \\
\bottomrule
\end{tabular}
\caption{Сравнение методов на $f_3$ при $\varepsilon=10^{-8}$}
\end{table}

\begin{figure}[H]
\centering
\includegraphics[width=\textwidth]{results/plots/func_3__comparison.png}
\caption{Зависимость числа итераций/вычислений от точности для $f_3$}
\end{figure}

\begin{figure}[H]
\centering
\includegraphics[width=\textwidth]{results/plots/func_3__interval_dynamics.png}
\caption{Динамика интервала неопределённости для $f_3$}
\end{figure}

Отдельно выполнено сравнение золотого сечения с/без разведки:
\begin{table}[H]
\centering
\small
\begin{tabular}{lrr}
\toprule
Сценарий & $x^*$ & $f(x^*)$ \\
\midrule
Без разведки (golden на $[-2,5]$) & 2.146776303 & $-0.5476338287$ \\
С разведкой (passive $\rightarrow$ golden) & 0.984052388 & $-0.7122701662$ \\
\bottomrule
\end{tabular}
\caption{Влияние разведки на результат для $f_3$}
\end{table}

Разведка пассивным поиском (1000 точек) дала стартовый интервал
\[
[0.484985,\; 1.484985],
\]
что позволило активному методу выйти к более глубокому минимуму.

\textbf{Краткий вывод:} Из графика видно, что разные методы сходятся к различным локальным минимумам. Это связано с тем, что функция является многомодальной, и результат зависит от начального интервала и устройства самого метода оптимизации.  \\
Методы, не использующие разведку, могут сходиться к менее глубоким минимумам.

\section{Сравнение с библиотечным методом Брента}
\begin{itemize}
    \item На гладкой унимодальной $f_1$ библиотечный Брент и метод парабол дают лучшие затраты.
    \item На функции с плато $f_{2.1}$ Брент корректен, но выбор точки минимума не уникален (плато).
    \item На асимметричной $f_{2.2}$ Брент и интервальные методы устойчиво находят минимум; метод парабол заметно хуже по числу итераций и вычислений.
    \item На многомодальной $f_3$ Брент без разведки может прийти в локальный минимум, поэтому разведка критична.
\end{itemize}

\section{Выводы}
\begin{enumerate}
    \item Все требуемые методы реализованы и протестированы на заданных типах функций.
    \item Для унимодальных функций построены таблицы зависимости от $\varepsilon$, графики сравнения и графики динамики интервала.
    \item Пассивный поиск полезен как разведочный инструмент, но дорог по числу вычислений.
    \item На сложных функциях надёжнее всего работают интервальные методы и Брент; параболы чувствительны к форме функции.
    \item Для многомодальной задачи комбинация ``разведка + активный метод'' даёт существенно лучший минимум, чем запуск активного метода без разведки.
\end{enumerate}

\end{document}
